\section{Synthesis, verification and the remaining questions}\label{art:verification}

The mixed benchmark $u^2+|v|$ has a whole-sequence convergence answer within
an arbitrary-dimensional centered-family theorem. Finite polyhedral geometry
supplies a local inward model, while a moving field handles the
three-quadratic example's curved switches. In each case comparison energies
use the principal-minor hierarchy to detect a remaining objective gap.

\paragraph{Why the counterexamples evade these hypotheses.}
The distinctions can be read from the constructed objectives.
The homogeneous seed satisfies $G(t(-1,16))=-t$, so it is unbounded below;
its coercive cap contains the affine term $v-17$ and is not homogeneous.
Neither belongs to the coercive homogeneous class of
Theorem~\ref{art:homogeneous}. In the local embeddings, wrong vertices recur
at every scale around the same minimizer. They cannot be removed by taking
a smaller neighborhood, and the resulting switching geometry is not resolved
by a finite local polyhedral-cell partition as in
Theorem~\ref{art:polyhedral}. Finitely many smooth formulas do not by themselves
supply that partition.

There is also a direct test of the planar hypotheses. In
\eqref{art:upper_objective}, all four branches are active at zero, with
indexed gradients $0,0,0,e_2$, where $e_2=(0,1)^T$.
Equation~\eqref{art:lower_objective} gives the same pattern before whitening:
the common quadratic has zero gradient at zero, and $\psi_b'(0)=0$.
After the one fixed change $x=Az$, the gradients are $0,0,0,A^Te_2$,
with $A^Te_2\ne0$. Thus the repeated zero gradients violate affine
independence, and a combination with positive weight $\lambda_3$ on the
fourth branch equals $\lambda_3e_2\ne0$ (or
$\lambda_3A^Te_2\ne0$ after whitening). The strictly positive optimal
weights in Corollary~\ref{art:planar-max} therefore fail as well.
The lower-curvature branches also lose second differentiability at zero.
Positive strong convexity alone supplies neither the missing regularity nor
this nondegeneracy.

The secant calculation explains how curvature and the failure mechanism
coexist. For $Q_\mu(x)=\mu\|x\|^2/2+G(x)$, suppose the differentiable
endpoints $x$ and $x+s$ have active slopes $g_i$ and $g_j$. Then exactly
\begin{equation}\label{art:jump-secant}
 y=\bigl(\mu(x+s)+g_j\bigr)-\bigl(\mu x+g_i\bigr)
   =\mu s+(g_j-g_i).
\end{equation}
Along the certified perturbed orbit, $s\to0$ while the three distinct
slopes continue to cycle. The positive-curvature contribution therefore
vanishes with the step, while the branch jumps persist. The stable-manifold
and strict-query arguments establish that this decomposition is realized
by the specified algorithm; a branch change alone does not imply failure.

These comparisons explain sufficient hypotheses separating the proved
positive classes from the examples. They do not make affine independence,
positive multipliers or a comparison certificate necessary for convergence,
and they do not turn every nonoptimal affine-critical point or accumulating
singularity into a counterexample. The broader local question for finitely
many uniformly strongly convex $C^2$ branches remains open here.
The examples concern the specified weak-Wolfe implementation and exceptional
deterministic starts; they do not settle Lewis--Overton's almost-sure
challenge under its additional hypotheses~\cite[Challenge~7.1]{LO}.

\paragraph{Formal verification and reproducibility.}
The complete companion supplies pinned Lean sources, exact theorem contracts
and their correspondence to the manuscripts, compiler and dependency audits,
and versioned source/build records. The proofs and rational certificates
needed for the article's main results are also contained here.
Formal checking and author-side correspondence review have distinct roles;
neither certifies publication priority or floating-point robustness.

\paragraph{AI assistance.}
Large language models assisted theorem development, proof reconstruction,
formalization, literature searches, software and writing. The companion
documents this substantive research assistance and the automated author-side
reviews.
